% Generated draft: requires mathematical and source review.
\documentclass[11pt]{article}
\usepackage[margin=25mm]{geometry}
\usepackage{fontspec}
\setmainfont{Latin Modern Roman}
\usepackage{amsmath,amssymb,mathtools,unicode-math}
\setmathfont{Latin Modern Math}
\usepackage{enumitem,needspace,xurl,hyperref}
\setlength{\parindent}{0pt}
\setlength{\parskip}{5pt}
\setlength{\emergencystretch}{3em}
\newcommand{\N}{\mathbb{N}}
\newcommand{\Z}{\mathbb{Z}}
\newcommand{\Q}{\mathbb{Q}}
\newcommand{\R}{\mathbb{R}}
\newcommand{\C}{\mathbb{C}}
\newcommand{\F}{\mathbb{F}}
\newcommand{\A}{\mathbb{A}}
\newcommand{\PP}{\mathbb{P}}
\newcommand{\E}{\mathbb{E}}
\newcommand{\eps}{\varepsilon}
\newcommand{\dd}{\,\mathrm{d}}
\begin{document}
\section*{UnsolvedMath ID: 800141}
\textbf{Source problem code:} catalogue entry
\subsection{Definitions and mathematical statement}
\textbf{Quadratic-form prime counts (Stanley Yao Xiao).}
Let
\[
f(U,V)=AU^2+BUV+CV^2\in\Z[U,V],\qquad \gcd(A,B,C)=1,
\]
be positive definite, meaning that \(f(u,v)>0\) for every \((u,v)\in\R^2\setminus\{(0,0)\}\). Write \(\operatorname{Disc}(f)=B^2-4AC\). Impose the source's condition
\[
f(T,1)\not\equiv T(T+1)\pmod 2.
\]
For real \(x\ge2\), define
\[
\pi_f(x)=\#\{p\le x:p\text{ is prime and }p=f(a,b^2)\text{ for some }(a,b)\in\Z^2\}.
\]
Each prime is counted once, irrespective of the number of representations. The supplied source attributes to reference [X] the asymptotic
\[
\pi_f(x)\sim C(\operatorname{Disc}(f))\frac{x^{3/4}}{\log x}
\qquad(x\to\infty),
\]
where the constant is positive and depends only on the discriminant. Here \(h(x)\sim k(x)\) means \(h(x)/k(x)\to1\). This attributed result is retained as supplied background, without independent verification.

The question asks: for two distinct forms \(f,g\) satisfying these conditions and having equal discriminants, how do \(\pi_f(x)\) and \(\pi_g(x)\) compare as \(x\to\infty\)? The source does not specify a comparison statistic, error scale, eventual inequality, or density of an inequality set. It also does not specify whether distinct forms must be inequivalent under a change of variables. Thus the requested comparison beyond the supplied common leading asymptotic is not a fully specified proposition.

The source additionally proposes an analogous comparison for indefinite binary quadratic forms, with the representing integers \(a,b\) restricted to a suitable box depending on \(x\). No box, size restrictions, or corresponding counting function is supplied, so this extension cannot be stated precisely from the record.

For primitive binary cubic forms \(F\in\Z[U,V]\), where primitive means that the coefficients have greatest common divisor \(1\), the source imposes
\[
F(U,V)\not\equiv UV(U+V)\pmod2.
\]
It attributes a similar asymptotic formula to Heath-Brown and Moroz [HBM] and asks for distinct such forms \(F,G\) with \(\pi_F(x)\sim\pi_G(x)\), followed by a comparison of their counting functions. The source does not define \(\pi_F\), specify its representation domain or any bounding region, or state the cubic asymptotic. These missing data preclude a precise cubic formulation; the quadratic definition is not silently extended to this setting.

\textbf{Weighted Chebyshev bias (Shin-ya Koyama).}
For \(r\in\{1,3\}\) and real \(x\ge2\), define
\[
\pi(x;4,r)=\#\{p\le x:p\text{ is prime},\ p\equiv r\pmod4\},
\]
and put
\[
W(x)=\sum_{\substack{p\le x\\p\equiv3\!\!\!\pmod4}}\frac1{\sqrt p}
-\sum_{\substack{p\le x\\p\equiv1\!\!\!\pmod4}}\frac1{\sqrt p}.
\]
The proposed weighted formulation, attributed in the source to [AK], is the property
\[
W(x)\sim\frac12\log\log x\qquad(x\to\infty).
\tag{W}
\]
Let
\[
S=\{x\ge2:\pi(x;4,3)>\pi(x;4,1)\}.
\]
Its logarithmic density, when the limit exists, is
\[
\delta_{\log}(S)=\lim_{X\to\infty}\frac1{\log X}\int_2^X\mathbf1_S(t)\,\frac{\mathrm dt}{t},
\]
where \(\mathbf1_S\) is the indicator of \(S\). The question asks what logical implications or other relationship connect property \((\mathrm W)\) to the conventional description that \(\delta_{\log}(S)\) is very close to \(1\). No numerical meaning of ``very close,'' hypotheses for the claimed weighted asymptotic, or hypotheses ensuring existence of the density are supplied. Consequently, the record does not specify two precise propositions between which a particular implication can be asserted.

\textbf{Incomplete prime-race fragment (Greg Martin).}
The remaining text attributes a question to Carl Pomerance and reference [Mar, Section 3], then introduces distinct reduced residue classes \(a,b\) modulo \(q\). Reduced means \(\gcd(a,q)=\gcd(b,q)=1\), and distinct means \(a\not\equiv b\pmod q\). The source ends at this point: it supplies neither the question nor an explicit range for \(q\). No mathematical target for this fragment can be recovered without additional source text.
\subsection{Short English statement}
The record contains a comparison question for primes represented by quadratic forms, incompletely specified indefinite and cubic extensions, a question relating weighted Chebyshev bias to logarithmic density, and an unfinished prime-race question. The supplied leading asymptotic and reference attributions are retained, while missing comparison criteria, hypotheses, definitions, and source text remain explicitly unresolved.
\subsection{Sources}
\begin{itemize}
\item Local UnsolvedMath registry, permanent ID \(800141\), exact source problem code \texttt{catalogue entry}, title ``Comparative Prime Number Theory --- Problem 10,'' in the AMR Open Problem Lists collection. The supplied registry labels the record open; this is a supplied status label, not an independently verified mathematical assessment.
\item Hamieh, Kadiri, Martin, and Ng, \emph{Comparative Prime Number Theory Problem List} (2024), source item Problem (\#10), \url{https://arxiv.org/abs/2407.03530}. The record reports access on 2026-10-16 and extraction by \texttt{arxiv\_latex}; these are supplied provenance details, not a claim of independent source access or verification.
\item The quadratic-form question and its extensions are attributed in the supplied excerpt to Stanley Yao Xiao, University of Northern British Columbia. Reference [X] is described there as Xiao's paper containing the stated main theorem, but no bibliographic identification is supplied.
\item The cubic-form asymptotic is attributed in the supplied excerpt to Heath-Brown and Moroz, reference [HBM]. No title, publication details, or separate link is supplied.
\item The weighted Chebyshev-bias question is attributed to Shin-ya Koyama, Toyo University. Its cited reference [AK] is not bibliographically identified in the supplied record.
\item The final fragment is attributed to Greg Martin, University of British Columbia; it states that Carl Pomerance asked the question and cites [Mar, Section 3]. The reference is not bibliographically identified, and the question itself is absent from the supplied excerpt.
\end{itemize}
\end{document}
